\chapter{AI Methodology and Model Development}
\label{ch:ai_methodology}

\section{Introduction}
\label{sec:ai_intro}
The central research inquiry of this senior project centers on comparing distinct algorithmic paradigms for satellite change detection in hyper-arid environments. This chapter provides a rigorous mathematical and structural exposition of all artificial intelligence models developed, calibrated, and evaluated within TerraShift AI:
\begin{enumerate}
    \item \textbf{Deep Siamese Convolutional U-Net} with difference-attention skip connections (\texttt{siamese\_untuned} and \texttt{siamese\_tuned\_filtered}).
    \item \textbf{AlphaEarth Model A (\texttt{alphaearth\_cosine}):} Unsupervised 64-dimensional latent cosine dissimilarity ($\tau=0.12$).
    \item \textbf{AlphaEarth Model B (\texttt{alphaearth\_dynamic\_world}):} Foundation embedding divergence fused with Dynamic World LULC probability transitions.
    \item \textbf{Standalone Dynamic World V1 (\texttt{dynamic\_world\_lulc}):} Google Dynamic World 9-class built-up and agricultural probability transition detector.
    \item \textbf{TerraShift v6.0 Hybrid (\texttt{alphaearth\_hybrid}, Model D):} Multi-modal fusion of Siamese U-Net spatial boundaries, AlphaEarth 64-D embeddings, and Sentinel-2 L2A directional physical spectral guards.
    \item \textbf{Directional Physical Spectral Guards} formulated from electromagnetic reflectance physics ($\Delta\text{NDBI}, \Delta\text{NDVI}$, and albedo drift bounds).
    \item \textbf{Exact WGS84 Bounding-Box Clipping \& Connected Component Polygonization}.
\end{enumerate}

\section{Siamese U-Net Deep Learning Architecture}
\label{sec:ai_siamese}
The core deep neural network in TerraShift AI is a modified **Fully Convolutional Siamese U-Net** (Figure \ref{fig:siamese_network_architecture}), built upon the foundational principles of \citet{daudt2018fully} and \citet{ronneberger2015unet}.

\begin{figure}[htbp]
\centering
\fbox{
\begin{minipage}{0.92\textwidth}
\centering
\vspace*{0.3cm}
\textbf{\large TerraShift Siamese U-Net Architectural Topology}\\[0.3cm]
\begin{tabular}{|c|c|c|}
\hline
\textbf{Branch 1: Pre-Change ($I_1$)} & \textbf{Shared Weights ($W_\theta$)} & \textbf{Branch 2: Post-Change ($I_2$)} \\
$I_1 \in \R^{H \times W \times 3}$ & $\longleftrightarrow$ & $I_2 \in \R^{H \times W \times 3}$ \\
\hline
$\downarrow$ Conv Block 1 ($64$) & Shared Conv Filters & $\downarrow$ Conv Block 1 ($64$) \\
$\downarrow$ MaxPool ($H/2, W/2$) & $\longleftrightarrow$ & $\downarrow$ MaxPool ($H/2, W/2$) \\
$\downarrow$ Conv Block 2 ($128$) & Shared Conv Filters & $\downarrow$ Conv Block 2 ($128$) \\
$\downarrow$ MaxPool ($H/4, W/4$) & $\longleftrightarrow$ & $\downarrow$ MaxPool ($H/4, W/4$) \\
$\downarrow$ Conv Block 3 ($256$) & Shared Conv Filters & $\downarrow$ Conv Block 3 ($256$) \\
$\downarrow$ MaxPool ($H/8, W/8$) & $\longleftrightarrow$ & $\downarrow$ MaxPool ($H/8, W/8$) \\
$\downarrow$ Conv Block 4 ($512$) & Shared Conv Filters & $\downarrow$ Conv Block 4 ($512$) \\
\hline
\multicolumn{3}{|c|}{$\Downarrow$ \textbf{Absolute Difference Latent Bottleneck: } $\mathbf{d}^{(4)} = |\mathbf{z}_1^{(4)} - \mathbf{z}_2^{(4)}|$} \\
\hline
\multicolumn{3}{|c|}{$\Uparrow$ \textbf{Decoder Level 3:} UpConv ($256$) $\oplus$ Difference Skip $\mathbf{d}^{(3)} = |\mathbf{z}_1^{(3)} - \mathbf{z}_2^{(3)}|$} \\
\multicolumn{3}{|c|}{$\Uparrow$ \textbf{Decoder Level 2:} UpConv ($128$) $\oplus$ Difference Skip $\mathbf{d}^{(2)} = |\mathbf{z}_1^{(2)} - \mathbf{z}_2^{(2)}|$} \\
\multicolumn{3}{|c|}{$\Uparrow$ \textbf{Decoder Level 1:} UpConv ($64$) $\oplus$ Difference Skip $\mathbf{d}^{(1)} = |\mathbf{z}_1^{(1)} - \mathbf{z}_2^{(1)}|$} \\
\multicolumn{3}{|c|}{$\Downarrow$ \textbf{Output Layer: } $1 \times 1$ Conv $\to$ Sigmoid $\to \hat{P}_{\text{change}} \in [0, 1]^{H \times W}$} \\
\hline
\end{tabular}
\vspace*{0.3cm}
\end{minipage}
}
\caption{Siamese U-Net architecture with multi-scale absolute difference skip connections.}
\label{fig:siamese_network_architecture}
\end{figure}

\subsection{Mathematical Formulation}
Let $I_1, I_2 \in \R^{H \times W \times C}$ denote co-registered bitemporal satellite optical patches. The parameter-shared encoder $f_\theta$ processes each image independently across four hierarchical resolution stages $l \in \{1, 2, 3, 4\}$:
\begin{equation}
    \mathbf{z}_1^{(l)} = f_\theta^{(l)}(I_1), \quad \mathbf{z}_2^{(l)} = f_\theta^{(l)}(I_2)
    \label{eq:encoder_stages}
\end{equation}
At each encoder stage, multi-scale feature difference maps $\mathbf{d}^{(l)}$ are computed using the element-wise absolute difference:
\begin{equation}
    \mathbf{d}^{(l)} = \abs{\mathbf{z}_1^{(l)} - \mathbf{z}_2^{(l)}} \in \R^{H_l \times W_l \times C_l}
    \label{eq:difference_skip}
\end{equation}
The absolute difference operation enforces exact temporal symmetry:
\begin{equation}
    \abs{f_\theta^{(l)}(I_1) - f_\theta^{(l)}(I_2)} \equiv \abs{f_\theta^{(l)}(I_2) - f_\theta^{(l)}(I_1)}
    \label{eq:temporal_symmetry}
\end{equation}
guaranteeing that flipping the temporal order of observation images produces mathematically identical change activation magnitudes.

The decoder synthesizes features from the bottleneck upward using transposed convolutions, concatenating each upsampled feature map with its corresponding difference skip connection $\mathbf{d}^{(l)}$:
\begin{equation}
    \mathbf{y}^{(l-1)} = \text{Conv}\left( \left[ \text{UpConv}(\mathbf{y}^{(l)}) \,\|\, \mathbf{d}^{(l-1)} \right] \right)
    \label{eq:decoder_upsample}
\end{equation}
The final layer applies a $1 \times 1$ convolution followed by a sigmoid activation function $\sigma(z) = (1 + e^{-z})^{-1}$, generating a continuous change posterior probability map $\hat{P} \in [0, 1]^{H \times W}$:
\begin{equation}
    \hat{P}(x, y) = \sigma\left( \mathbf{w}_{\text{out}}^T \mathbf{y}^{(0)}(x, y) + b_{\text{out}} \right)
    \label{eq:sigmoid_output}
\end{equation}

\subsection{Training Loss and Severe Class Imbalance}
In operational Earth observation, change detection is plagued by **extreme class imbalance**: in a typical satellite tile, genuine change pixels comprise less than $1\%\text{ to }3\%$ of total surface area, while unchanged background pixels dominate $97\%\text{ to }99\%$ of the scene. Standard Binary Cross-Entropy (BCE) loss on uncalibrated networks leads to severe failure, as the optimizer trivially minimizes loss by predicting zero change everywhere.

To overcome this, TerraShift's Siamese network is trained using a compound loss function coupling Weighted Binary Cross-Entropy with Soft Generalized Dice Loss \citep{sudre2017generalised}:
\begin{equation}
    \mathcal{L}_{\text{total}} = \lambda_{\text{BCE}} \mathcal{L}_{\text{wBCE}} + \lambda_{\text{Dice}} \mathcal{L}_{\text{Dice}}
    \label{eq:compound_loss}
\end{equation}
The weighted BCE loss penalizes false negatives heavily:
\begin{equation}
    \mathcal{L}_{\text{wBCE}} = -\frac{1}{N} \sum_{i=1}^N \left[ w_1 y_i \log(\hat{p}_i) + (1 - y_i) \log(1 - \hat{p}_i) \right]
    \label{eq:weighted_bce}
\end{equation}
where $w_1 = (N - N_{\text{pos}}) / N_{\text{pos}}$ dynamically balances positive change pixels.

The Generalized Soft Dice loss optimizes region overlap directly:
\begin{equation}
    \mathcal{L}_{\text{Dice}} = 1.0 - \frac{2 \sum_{i=1}^N y_i \hat{p}_i + \epsilon}{\sum_{i=1}^N y_i + \sum_{i=1}^N \hat{p}_i + \epsilon}
    \label{eq:soft_dice}
\end{equation}
where $\epsilon = 10^{-6}$ provides numerical stability. Training hyperparameters are set to $\lambda_{\text{BCE}} = 0.5$, $\lambda_{\text{Dice}} = 0.5$, batch size of 16, AdamW optimizer with initial learning rate $\eta = 3 \times 10^{-4}$, and cosine annealing decay.

\section{Model A: Unsupervised AlphaEarth Cosine Dissimilarity (\texttt{alphaearth\_cosine})}
\label{sec:ai_model_a}
AlphaEarth Model A evaluates Google's 64-dimensional annual satellite embeddings in an entirely unsupervised metric-space formulation.

\subsection{Unit-Sphere Cosine Formulation}
For each pixel coordinate $(x, y)$, let $\mathbf{u}(x, y) \in \R^{64}$ denote the AlphaEarth embedding for year $T_1$ and $\mathbf{v}(x, y) \in \R^{64}$ denote the embedding for year $T_2$. Because AlphaEarth vectors are strictly unit-normalized by construction ($\norm{\mathbf{u}}_2 = 1.0$, $\norm{\mathbf{v}}_2 = 1.0$), their inner product equals their cosine similarity:
\begin{equation}
    S_{\text{cosine}}(x, y) = \sum_{k=0}^{63} u_k(x, y) \cdot v_k(x, y)
    \label{eq:dot_prod_cosine}
\end{equation}
The cosine dissimilarity metric $D_{\text{cosine}} \in [0, 2]$ is computed as:
\begin{equation}
    D_{\text{cosine}}(x, y) = 1.0 - \text{clip}\left( S_{\text{cosine}}(x, y), -1.0, 1.0 \right)
    \label{eq:cosine_dist_eval}
\end{equation}
A continuous normalized change score $S_{\text{change}}(x, y) \in [0, 1]$ is generated by scaling over the active dissimilarity range:
\begin{equation}
    S_{\text{change}}(x, y) = \text{clip}\left( \frac{D_{\text{cosine}}(x, y)}{0.50}, 0.0, 1.0 \right)
    \label{eq:scaled_change_score}
\end{equation}
Applying the optimal screening threshold $\tau_{\text{dist}} = 0.12$ established by \citet{manikandan2026alphaearth} (corresponding to cosine similarity $< 0.88$) generates the candidate binary change mask:
\begin{equation}
    M_A(x, y) = \begin{cases} 
        1, & \text{if } D_{\text{cosine}}(x, y) \ge 0.12 \land V(x, y) = 1 \\ 
        0, & \text{otherwise} 
    \end{cases}
    \label{eq:mask_model_a}
\end{equation}

\section{Model B \& Standalone Dynamic World Transition Analysis (\texttt{alphaearth\_dynamic\_world} and \texttt{dynamic\_world\_lulc})}
\label{sec:ai_model_b}
To evaluate whether global land-cover probability models can solve arid change detection alone or in combination with AlphaEarth, we formulate both **Standalone Google Dynamic World V1 LULC Transition (\texttt{dynamic\_world\_lulc})** and **Model B (\texttt{alphaearth\_dynamic\_world})** \citep{brown2022dynamicworld}.

\subsection{Semantic Transition Verification}
Let $\mathbf{P}_1(x, y), \mathbf{P}_2(x, y) \in [0, 1]^9$ denote the Dynamic World 9-class probability vectors at dates $T_1$ and $T_2$. In Saudi development corridors, genuine anthropogenic changes are dominated by two primary class transitions:
\begin{enumerate}
    \item \textbf{Built-Up Infrastructure Expansion:} A positive surge in built-up land probability:
    \begin{equation}
        \Delta P_{\text{built}}(x, y) = P_2(x, y)[\text{built}] - P_1(x, y)[\text{built}] \ge +0.15
        \label{eq:dw_built_delta}
    \end{equation}
    \item \textbf{Agricultural Dynamic Shift:} A significant change in crop probability:
    \begin{equation}
        \abs{\Delta P_{\text{crops}}(x, y)} = \abs{P_2(x, y)[\text{crops}] - P_1(x, y)[\text{crops}]} \ge 0.20
        \label{eq:dw_crop_delta}
    \end{equation}
\end{enumerate}
Standalone Dynamic World (\texttt{dynamic\_world\_lulc}) triggers whenever $(\Delta P_{\text{built}} \ge 0.12 \lor \abs{\Delta P_{\text{crops}}} \ge 0.18) \land V(x,y)$, whereas Model B (\texttt{alphaearth\_dynamic\_world}) requires both latent vector divergence and semantic class transition:
\begin{equation}
    M_B(x, y) = \left( D_{\text{cosine}}(x, y) \ge 0.10 \right) \land \left( \Delta P_{\text{built}}(x, y) \ge 0.15 \lor \abs{\Delta P_{\text{crops}}(x, y)} \ge 0.20 \right) \land V(x, y)
    \label{eq:mask_model_b}
\end{equation}

\subsection{Methodological Exclusion of Unverified Supervised Classifiers}
While our codebase retains a 257-dimensional supervised classifier interface (\texttt{AlphaEarthDifferenceClassifier}) for future use with surveyed cadastral pixel masks, we strictly prohibit fitting a supervised classifier on coarse parcel bounding boxes or synthetic calibration distributions during official benchmarks. Calling \texttt{predict()} without an explicitly fitted ground-truth model raises a \texttt{RuntimeError}, verified by unit test \texttt{test\_alphaearth\_classifier\_requires\_explicit\_fit}.

\section{Model D: TerraShift v6.0 Hybrid (\texttt{alphaearth\_hybrid})}
\label{sec:ai_model_d}
Model D (\texttt{alphaearth\_hybrid}) represents the pinnacle of our architecture: a multi-modal fusion framework that combines the **Siamese U-Net spatial boundary prior**, **Google AlphaEarth V1 64-D semantic embeddings**, and **Sentinel-2 Level-2A directional physical spectral guards**.

\subsection{Multi-Modal Spatial-Semantic-Physical Consensus}
Standalone AlphaEarth cosine distance ($D_{\text{cosine}} \ge 0.12$) suffers from desert albedo and coastal sabkha moisture false alarms, whereas the Siamese U-Net coupled with L2A spectral indices occasionally retains small peripheral soil-disturbance false positives. Model D fuses all three complementary modalities:
\begin{enumerate}
    \item \textbf{Built Infrastructure \& Earthworks Consensus:} A pixel is confirmed as urban/industrial construction if it passes the tuned Siamese + L2A physical mask $M_{\text{Siam+Phys}}(x,y) = 1$ (requiring $\Delta\text{NDBI} \ge +0.025$ or brightness/SWIR compaction shift) **AND** exhibits real latent representation divergence in AlphaEarth ($D_{\text{cosine}}(x,y) \ge 0.065$) or high structural probability ($P_{\text{Siam}}(x,y) \ge 0.66$).
    \item \textbf{Agricultural Pivot Consensus:} A pixel is confirmed as agricultural canopy activation or fallowing when $\abs{\Delta\text{NDVI}(x,y)} \ge 0.12$ is corroborated by AlphaEarth latent divergence ($D_{\text{cosine}}(x,y) \ge 0.075$).
    \item \textbf{Strong Semantic-Physical Discovery:} High-magnitude AlphaEarth shifts ($D_{\text{cosine}}(x,y) \ge 0.16$) that simultaneously satisfy strict L2A physical verification ($\Delta\text{NDBI} \ge +0.045 \lor \abs{\Delta\text{NDVI}} \ge 0.14$) are also retained.
\end{enumerate}

\subsection{Composite Decision Rule}
The final binary change mask for TerraShift v6.0 Hybrid (\texttt{alphaearth\_hybrid}) is:
\begin{equation}
    M_{\text{Hybrid}}(x, y) = \Big[ \big( M_{\text{Siam+Phys}}(x, y) \land (D_{\text{cosine}} \ge 0.065 \lor P_{\text{Siam}} \ge 0.66) \big) \lor \big( D_{\text{cosine}} \ge 0.16 \land (\Delta\text{NDBI} \ge 0.045 \lor \abs{\Delta\text{NDVI}} \ge 0.14) \big) \Big] \land V(x, y)
    \label{eq:mask_model_d}
\end{equation}
This tri-modal formulation suppresses **91.2\%** of untuned false alarms ($114 \to 10$) and eliminates **100\%** of negative control zone false alarms ($0$ hits across all 4 corridors) while preserving **100.0\%** event recall ($19/19$).

\section{Connected Component Analysis, Exact AOI Clipping, and Minimum-Area Pruning}
\label{sec:ai_polygonization}
Raw pixel-level classification masks are converted into analysis-ready vector GIS layers through an automated topological polygonization pipeline (\texttt{services/ml/terrashift/geo/components.py}):

\begin{enumerate}
    \item \textbf{8-Connectivity Component Labeling:} Contiguous clusters of active change pixels are segmented using 2D 8-connectivity connected-component labeling:
    \begin{equation}
        \mathcal{C} = \{ C_1, C_2, \dots, C_K \}
        \label{eq:connected_components}
    \end{equation}
    \item \textbf{Geometric Boundary Extraction \& Exact AOI Clipping:} For each connected component $C_k$, exterior boundary coordinates are transformed into geographic lon/lat coordinates $(\lambda, \phi)$ in \texttt{EPSG:4326} and geometrically intersected with the exact corridor bounding box $\mathcal{B}_{\text{AOI}} = \text{box}(\lambda_{\min}, \phi_{\min}, \lambda_{\max}, \phi_{\max})$:
    \begin{equation}
        P_k^{\text{clipped}} = P_k \cap \mathcal{B}_{\text{AOI}}
        \label{eq:aoi_clipping}
    \end{equation}
    This guarantees that no vectorized polygon ever extends outside the target AOI or exceeds $100.0\%$ of the corridor surface area.
    \item \textbf{Geodetic WGS84 Surface Area Measurement:} True physical surface area in square meters ($A_k$) is integrated on the WGS84 ellipsoid via \texttt{pyproj.Geod(ellps="WGS84")}:
    \begin{equation}
        A_k = \abs{\text{Geod.geometry\_area\_perimeter}(P_k^{\text{clipped}})[0]}
        \label{eq:geod_area}
    \end{equation}
    \item \textbf{Minimum-Area Spatial Pruning:} Any polygon whose clipped WGS84 geodesic area is smaller than $A_{\min} = 3{,}500\text{ m}^2$ is pruned:
    \begin{equation}
        \mathcal{C}_{\text{pruned}} = \{ P_k^{\text{clipped}} \mid A_k \ge 3{,}500\text{ m}^2 \}
        \label{eq:area_pruning}
    \end{equation}
\end{enumerate}
\newpage
